% ID: SEL-D0080-OA-E01
% Tanda S080.
Sean $x$ e $y$ los kilogramos de manzanas A y B. El presupuesto, la capacidad de transporte y la no negatividad dan
\[
 0{,}60x+y\le1200,\qquad x+y\le1500,\qquad x,y\ge0.
\]
El beneficio por kilogramo es el precio de venta menos el de compra, de modo que maximizamos $Z=0{,}30x+0{,}35y$ en euros.

La región factible es un polígono cerrado y acotado. Sus vértices sobre los ejes son $(0,0)$, $(1500,0)$ y $(0,1200)$. La intersección de las rectas de presupuesto y capacidad se obtiene restando sus ecuaciones, $0{,}40x=300$, y da $(750,750)$.
\[
 \begin{array}{c|rrrr}
 (x,y)&(0,0)&(1500,0)&(750,750)&(0,1200)\\
 Z\text{ (euros)}&0&450&487{,}50&420
 \end{array}
\]
El máximo se alcanza comprando 750 kg de cada tipo y es de $\boxed{487{,}50}$ euros. Esta compra cuesta 1200 euros y ocupa exactamente los 1500 kg de capacidad.
